\section{Исходная программа}

Исходный код родительского процесса:

\begin{center}
\lstinputlisting[
    language=C,
    caption={Исходный код родительского процесса},
    captionpos=b
]{inc/parent.c}
\end{center}

Исходный код дочернего процесса:

\begin{center}
\lstinputlisting[
    language=C,
    caption={Исходный код дочернего процесса},
    captionpos=b
]{inc/child.c}
\end{center}

\subsection{Трассировка системных вызовов}

Для исследования системных вызовов программа была запущена с помощью
утилиты \texttt{strace}:

\begin{lstlisting}
strace -f -tt -s 200 -o strace.log ./parent
\end{lstlisting}

Параметр \texttt{-f} позволяет отслеживать не только родительский,
но и созданный им дочерний процесс.

Ключевой фрагмент трассировки создания канала и дочернего процесса
имеет следующий вид:

\begin{lstlisting}
pipe2([3, 4], 0) = 0
clone(child_stack=NULL, flags=CLONE_CHILD_CLEARTID|CLONE_CHILD_SETTID|SIGCHLD <unfinished ...>)
close(4) = 0
dup2(3, 0) = 0
close(3) = 0
execve("./child", ["child", "1 2.5"], 0x7ffdbe8cdfb8 /* 72 vars */ <unfinished ...>)
\end{lstlisting}

Вызов \texttt{pipe2()} создаёт канал с файловыми дескрипторами 3
и 4. Дескриптор 3 используется для чтения, а дескриптор 4 ---
для записи. В дочернем процессе вызов \texttt{dup2(3, 0)}
перенаправляет читающий конец канала на файловый дескриптор 0,
то есть на стандартный поток ввода.

Передача строки между процессами отражается в трассировке следующим
образом:

\begin{lstlisting}
write(4, "3.4 5.4 7.6\n", 12) = 12
read(0, "3.4 5.4 7.6\n", 4096) = 12
\end{lstlisting}

Первый вызов выполняется родительским процессом и записывает строку
в канал. Второй вызов выполняется дочерним процессом. Поскольку его
стандартный ввод был перенаправлен на читающий конец канала,
\texttt{read(0, ...)} получает данные, переданные родителем.

После завершения пользовательского ввода родитель получает EOF и
закрывает пишущий конец канала:

\begin{lstlisting}
<... read resumed>("", 1024) = 0
close(4) = 0
\end{lstlisting}

После закрытия последнего пишущего дескриптора дочерний процесс также
получает EOF:

\begin{lstlisting}
read(0, "", 4096) = 0
\end{lstlisting}

Родительский процесс ожидает завершения дочернего процесса. В трассировке
вызов \texttt{waitpid()} представлен системным вызовом \texttt{wait4()}.

\begin{lstlisting}
wait4(14059,  <unfinished ...>)
(<... wait4 resumed>[{WIFEXITED(s) && WEXITSTATUS(s) == 0}], 0, NULL) = 14059
--- SIGCHLD {si_signo=SIGCHLD, si_code=CLD_EXITED, si_pid=14059, si_uid=1000, si_status=0, si_utime=0, si_stime=0} ---
exit_group(0)     = ?
\end{lstlisting}

Перед завершением дочерний процесс записывает результаты в выходной файл,
закрывает его и завершается с кодом 0. После завершения дочернего процесса
родителю передаётся сигнал \texttt{SIGCHLD}, а ожидающий системный вызов
\texttt{wait4()} возвращает статус завершения дочернего процесса.